\section*{第 \ref{ch:b3:rna-regulation} 章——非编码 RNA 与转录后调控}

\begin{solution}{exo:b3:rna-regulation:1}
\omterm{def:b3:rna-regulation:ncrna}{miRNA}：以发夹结构编码于基因组中，\qty{22}{nt}，由 \omterm{def:b3:rna-regulation:mirna}{Drosha} 切割
随后由 \omterm{def:b3:rna-regulation:rnai}{Dicer} 加工，经种子区不完全配对，抑制
许多信使 RNA 的翻译与稳定性。\omterm{def:b3:rna-regulation:ncrna}{siRNA}：来源于长
双链 RNA（来自病毒、转基因或实验），\qty{21}{nt}，
依赖 \omterm{def:b3:rna-regulation:rnai}{Dicer} 切割，完全配对，并引导 \omterm{def:b3:rna-regulation:rnai}{Argonaute} 切割
靶标。\omterm{def:b3:rna-regulation:ncrna}{piRNA}：来源于 \omterm{def:b3:rna-regulation:ncrna}{piRNA} 簇的单链转录本，
\qtyrange{26}{31}{nt}，不依赖 \omterm{def:b3:rna-regulation:rnai}{Dicer}，在生殖系中装载于 PIWI 蛋白，
切割转座子转录本（乒乓循环），
并引导转座子 DNA 的\omterm{def:b3:chromatin-epigenetics:nucleosome}{染色质}沉默。
\end{solution}

\begin{solution}{exo:b3:rna-regulation:2}
Pol II 转录产生 pri-miRNA；\omterm{def:b3:rna-regulation:mirna}{Drosha}（与 DGCR8 共同作用）切下
发夹结构，于细胞核中完成；exportin-5 输出 pre-miRNA；\omterm{def:b3:rna-regulation:rnai}{Dicer} 切去
环区，留下 \qty{22}{nt} 双链体；其中一条链装载入
\omterm{def:b3:rna-regulation:rnai}{Argonaute}；种子区（第 2--8 位核苷酸）与 3$'$ UTR 位点配对；
\omterm{def:b3:rna-regulation:rnai}{Argonaute} 招募去腺苷酸酶与去帽装置，并阻断
翻译起始——信使 RNA 的翻译减少，降解加快。
\end{solution}

\begin{solution}{exo:b3:rna-regulation:3}
体外由噬菌体聚合酶合成的 RNA 含有双链
污染物（该聚合酶还会转录模板的另一条链，并在末端脱落），
因此“有义链”制剂中携带了真正的
触发物。法尔与梅洛分别纯化了每条链，证明单独任一条链
作用很小，随后有意将两者退火：双链 RNA
的效力至少高出百倍以上，每个细胞仅需几个分子
即可奏效。
\end{solution}

\begin{solution}{exo:b3:rna-regulation:4}
缺铁时：IRP 结合 IRE。铁蛋白：翻译起始
受阻，不合成储存蛋白（翻译调控）。
转铁蛋白受体：其信使 RNA 不受核酸酶降解，
因而积累，受体合成增加，铁的摄入增多（
信使 RNA 稳定性调控）。
\end{solution}

\begin{solution}{exo:b3:rna-regulation:5}
$\gamma = \ln 2/\qty{120}{min} = \qty{0.0058}{min^{-1}}$; $m^{*} =
5/0.0058 \approx 870$ 个分子。降解速率加倍：$m^{*} \approx 430$，
半衰期为 $\qty{120}{min}/2 = \qty{60}{min}$。
\end{solution}

\begin{solution}{exo:b3:rna-regulation:6}
3$'$ UTR 序列总长 $2\times 10^{7}$ nt；偶然匹配数为 $2\times
10^{7}/4^{7} = 2\times 10^{7}/\num{16384} \approx \num{1200}$ 个位点。
真正的靶标可依据以下特征区分：该位点在物种间保守，
所处环境有利（富含 AU 的邻近区域，位点远离
终止密码子，\omterm{def:b3:rna-regulation:ncrna}{miRNA} 的 3$'$ 端有额外配对），
\omterm{def:b3:rna-regulation:ncrna}{miRNA} 与靶标在同一细胞中共表达，并通过实验验证：
存在 \omterm{def:b3:rna-regulation:ncrna}{miRNA} 时信使 RNA 水平下降，
位点发生突变后则不再下降。
\end{solution}

\begin{solution}{exo:b3:rna-regulation:7}
(a) XX，\emph{tra} 无效突变：SXL 合成而 TRA 不合成，\emph{dsx} 采用
默认雄性剪接方式——体细胞发育为雄性（不育）。(b) 组成型表达 TRA 的 XY 个体：
TRA-2 在两种性别中均存在，因此 DSX 以
雌性形式合成——体细胞发育为雌性（一种不育的
“假雌性”，生殖系则随其他信号而定）。(c) 无
\emph{dsx}：两种 DSX 形式均不存在，因此两种终末
分化程序均未启动——形成间性个体，兼具混合或不完整的
性状。
\end{solution}

\begin{solution}{exo:b3:rna-regulation:8}
稳态下，信使 RNA 及由其决定的蛋白质量
与信使 RNA 半衰期成正比，故蛋白质多 $\qty{180}{min}/\qty{15}{min} = 12$
倍。该蛋白是一种正常的生长信号，本应是短暂的，
即每次刺激后出现的一次脉冲；若合成量高出十二倍且持续存在，
则它会在无刺激的情况下驱动细胞增殖。剂量与
时机，而非序列，使其具有致癌性。
\end{solution}

\begin{solution}{exo:b3:rna-regulation:9}
(a) 实验室母本、野生父本：卵中不携带针对
P 元件的 \omterm{def:b3:rna-regulation:ncrna}{piRNA}（母本谱系从未接触过它们），父本的 P 元件
在子代生殖系中自由转座——不育
（\omterm{def:b3:rna-regulation:pirna}{杂种不育}）。(b) 野生母本：其卵中含有 P 元件 \omterm{def:b3:rna-regulation:ncrna}{piRNA}，
使子代中的 P 元件沉默，可育。这种不对称性源于
母本向卵中沉积 \omterm{def:b3:rna-regulation:ncrna}{piRNA}，这是 RNA 的遗传，而非基因的遗传。
\end{solution}

\begin{solution}{exo:b3:rna-regulation:10}
蛋白质在其自身寿命 $\tau_{p}$ 内对其信使 RNA 进行积分。
信使 RNA 的寿命为 $\tau_{m} = 1/(\gamma + \gamma' R)$；在 $\tau_{p}$
期间，蛋白质因此约经历 $\tau_{p}/\tau_{m}$ 个相互独立的
信使 RNA 寿命，各为一个随机的生灭事件。相对
涨落在对 $N$ 个独立事件取平均时，按
$1/\sqrt{N}$ 下降，因此在微 RNA 调控下较短的 $\tau_{m}$
——在平均信使 RNA 量相同的前提下（这需要按比例提高 $k$）——
使蛋白质更平滑。细胞为微 RNA 及额外的转录
付出代价，以换取精确性与速度：由快速合成与快速降解
的平衡所确定的水平，噪声更小，且比由缓慢合成与缓慢降解
所确定的同一水平更易于重新调整。
\end{solution}

\begin{solution}{exo:b3:rna-regulation:11}
线虫：RNA 依赖性 RNA 聚合酶直接以靶标
信使 RNA 本身为模板合成新的 \omterm{def:b3:rna-regulation:ncrna}{siRNA}，只要靶标仍被转录，
且能经受分裂中的稀释，信号便持续更新；借助细胞间的
双链 RNA 通道，信号还可扩散至其他组织，并在数代内
进入生殖系。哺乳动物中则没有扩增，因此
\omterm{def:b3:rna-regulation:ncrna}{siRNA} 在每次分裂时被消耗和稀释，仅在接收它们的
细胞中起作用，并在分裂细胞中数日内消退。
由此可见：\omterm{def:b3:rna-regulation:ncrna}{siRNA} 药物须递送至每一个靶细胞，
效果最好的是不分裂的细胞（肝细胞，因此有肝脏
药物），并需要重复给药。
\end{solution}

\begin{solution}{exo:b3:rna-regulation:12}
(a) 功能性 RNA：转录产物合成后将其破坏（\omterm{def:b3:rna-regulation:ncrna}{siRNA}、
反义寡核苷酸）会产生表型，而在基因组其他位置以转基因
方式表达该 RNA 可使表型恢复。(b) 功能性转录：插入提前的
多聚腺苷酸化信号以终止转录（或删除启动子）会产生表型，
而转录后破坏 RNA 不会，异位转基因亦不能使其恢复——
起作用的是转录这一行为（或 DNA 元件），与许多
增强子相关 RNA 的情形一致。(c) 噪声：两种操作均
不改变任何结果，该 RNA 在每个细胞中的拷贝数
不足一个，
其序列与表达也不保守。
\end{solution}

\begin{solution}{pb:b3:rna-regulation:1}
\textbf{1.} $\gamma = \ln 2/40 = \qty{0.0173}{min^{-1}}$；$m^{*} =
3/0.0173 \approx 170$ 个信使。
\textbf{2.} $\delta_{p} = \ln 2/180 = \qty{0.00385}{min^{-1}}$；$P^{*} =
\beta m^{*}/\delta_{p} = 2\times 173/0.00385 \approx 9.0\times 10^{4}$
个蛋白质。
\textbf{3.} 降解速率 $4\gamma = \qty{0.069}{min^{-1}}$；$m^{*} = 3/0.069
\approx 43$；时间常数 $1/0.069 = \qty{14}{min}$。
\textbf{4.} $\beta = 1$：$P^{*} = 43/0.00385 \approx 1.1\times 10^{4}$；
阻遏 $8$ 倍（$4$ 来自降解，$2$ 来自翻译）。
\textbf{5.} 信使的半衰期 $\ln 2/(4\gamma) = \qty{10}{min}$。该
蛋白质随后以自身半衰期 \qty{3}{h} 衰减：
蛋白质的降解限制了开关。
\textbf{6.} 是的：信使仅降至 $1/4$（仍易于
检测），而蛋白质降至 $1/8$，并随旧蛋白
的更替持续下降；早期观察到蛋白质消失而
“mRNA 仍存在”，主要反映了阻遏中
翻译的那一半。
\textbf{7.} 每个卵中的分子数为 $10^{6}/250 = 4000$。
\textbf{8.} 孵化时每个细胞含 $4000/550 \approx 7$ 个。$4000/2^{n} < 1$
即自受精卵起 $n \ge 12$ 次分裂——约为
孵化后的三次分裂。
\textbf{9.} 该聚合酶以靶标信使为模板合成新的
双链 RNA，并切割成次级 \omterm{def:b3:rna-regulation:ncrna}{siRNA}：触发物由靶标本身
再生，因而能耐受稀释并扩散（双链 RNA 通道可在细胞间
传递双链 RNA，并进入生殖系）。信号消退，是因为扩增
同时需要靶标与初级触发物；注入的 RNA 消失后，次级
\omterm{def:b3:rna-regulation:ncrna}{siRNA} 池逐代稀释，
生殖系随之复位。
\textbf{10.} $5000/2^{n} < 100$：$2^{n} > 50$，$n = 6$ 天（剩余 $78$
个拷贝）。
\textbf{11.} 不分裂的细胞不会稀释 \omterm{def:b3:rna-regulation:rnai}{RISC}；限制因素是
\omterm{def:b3:rna-regulation:ncrna}{siRNA} 的化学寿命（核酸酶降解，可由链的化学修饰
减缓），以及向细胞质供给的
内体储库的缓慢渗漏。
\textbf{12.} $2\times 10^{7}/\num{16384} \approx \num{1200}$ 次
偶然匹配。遗传学确定：\emph{lin-14} 缺失产生的表型与
\emph{lin-4} 缺失相反，\emph{lin-14} 功能获得型等位基因
恰为 3$'$ UTR 位点的缺失，而七个位点
均保守。
\textbf{13.} 水平 $\propto$ 半衰期：$50/10 = 5$ 倍上升。
\textbf{14.} 铁蛋白合成下降 $1/0.05 = 20$ 倍。
\textbf{15.} 摄取与储存上升 $5\times 20 = 100$ 倍，由富铁
转为贫铁。
\textbf{16.} 只有胞质中有可用的铁时，铁硫簇才会组装，
因此蛋白质的构象直接读出了铁的
浓度，无需受体、激素或转录参与——类似于
\omterm{prop:b3:rna-regulation:translation}{核糖开关}，即代谢物被其所调控的分子感知。
\textbf{17.} 铁蛋白的翻译与铁无关：血清铁蛋白
极高，血清铁与储备量正常（无铁过载）。在晶状体中，
晶状体无更替，铁蛋白累积数年并结晶，
即早期白内障。
\textbf{18.} 发夹结构只是结合位点；效应取决于
结合的蛋白质物理上阻挡的对象：
5$'$ 端的扫描核糖体，3$'$ 端的核酸酶。
\textbf{19.} 两种化学机制覆盖两种条件：IRP1 的铁硫
簇也会被氧化剂与低氧破坏，因此 IRP1 响应
氧化还原状态，而 IRP2 的铁依赖性降解
直接响应铁本身，且在生理氧范围内均有效；冗余使
调控子稳健，两者报告的
是不同的信息。
\textbf{20.} $12\times 48\times 33\times 2 = \num{38016}$。$1 - (1 -
50/\num{38016})^{50} = 1 - 0.99868^{50} \approx 1 - e^{-0.066} =
0.064$：约 \qty{6}{\%}。
\textbf{21.} 相同的同工型通过同嗜性结合并相互排斥：
同一神经元的分支彼此回避、铺展开，而
不同神经元的分支几乎不共享同工型，可能相互交叉——
每个细胞各有其专属身份。
\textbf{22.} $2^{n}$；$\num{1024}$（$n = 10$ 时）。实际数目更少：
外显子是否被纳入由细胞类型专一的因子决定，并
在外显子之间相互协调；许多组合会改变阅读框，
并被无义介导的 mRNA 降解破坏，而且各外显子
并非彼此独立。
\textbf{23.} \emph{Sxl} 还控制\omterm{def:b3:chromatin-epigenetics:x-inactivation}{剂量补偿}：在雌性中
SXL 阻断 \emph{msl-2} 的翻译，因此两条 X 染色体
不发生超转录；若无 SXL，XX 胚胎会在两条 X 染色体上组装
MSL 复合体，使其产出加倍——对所有 X 连锁基因造成
致死性过量表达，甚至早于性别问题的出现。
\textbf{24.} 早期启动子仅在 X:A 信号存在时启动，
它产生的 SXL 促使来自维持启动子的转录本
进行有效剪接；维持启动子在两性中均有活性，因此一旦
SXL 存在，它便持续自我合成，不再需要该信号
——与\omterm{def:b3:chromatin-epigenetics:nucleosome}{染色质}标记中相同的、读取者招募书写者的
自我维持逻辑，此处则由蛋白质与一次剪接实现。
\textbf{25.} LIN-14 蛋白被阻遏 $8$ 倍；信使开关的
半衰期 \qty{10}{min}；未经扩增的 RNA 每个细胞的分子数
经 $12$ 次分裂后降至不足一个。
\end{solution}
